\bookchapter{第二十七章：爸媽都能輕鬆用——適老化手機完全指南}

\chapterintro{講真，寫呢一章嘅時候，我諗起我爸媽。我媽每次手機彈出一個「系統更新」，佢都唔敢點，驚整壞手機。我爸嘅微信字體比螞蟻仲細，但佢唔知點樣調大，就咁矇住眼睇咗兩年。呢個唔係佢哋「蠢」，係手機預設將所有人都當成二十幾歲嘅後生仔。呢一章，我用最直白嘅話，將長輩哋手機使用中最頭痛嘅問題一個個解開。}

\booksection{零、先將心態擺正：呢個唔係「老人用唔嚟」，係「手機冇做好」}

喺任何設定之前，有一個觀念要先講清楚：\textbf{而家嘅智能手機介面，由設計之初就冇點考慮老年用戶}。彈窗細到睇唔清、操作手勢需要精準觸碰、圖標生得都差唔多——呢啲唔係用戶嘅問題，係產品嘅問題。所以帶長輩用手機，第一件事唔係教，係幫佢哋\textbf{將環境調好}。環境調好咗，大半嘅問題就消失咗。

帶長輩用手機有三條鐵律：(1)\textbf{唔好催}——你用咗十年嘅智能機，佢哋可能先至用咗三個月。(2)\textbf{用類比}——「Wi-Fi就好似屋企嘅水管，開咗水龍頭先有水」，呢種比喻比專業術語好理解一萬倍。(3)\textbf{寫低}——口頭教一遍佢哋可能唔記得，你攞張紙將關鍵步驟寫低貼喺佢手機殼背面，比咩都管用。

\booksection{一、第一步：打開「簡易模式」——手機變清爽喇}

而家嘅手機基本都內置咗一個「簡易模式」（有啲叫「關懷模式」「老年模式」「長輩模式」），打開之後，成部手機都會大變樣：

\begin{itemize}
  \item \textbf{字體同圖標變得好大}——再亦唔使戴老花鏡睇手機喇。
  \item \textbf{桌面圖標大幅精簡}——只保留電話、短訊、相機、相冊、微信等最常用嘅幾個App，其他全部收埋。
  \item \textbf{操作變簡單喇}——好多花入面胡哨嘅功能（負一屏、智能助手、浮窗等）預設關閉，介面主打一個清爽。
  \item \textbf{聲音變大喇}——來電鈴聲、通知音量、通話聲音都自動調到最大。
\end{itemize}

\booksubsection{各品牌簡易模式開啟方法（直接跟住操作就得）：}

\begin{center}
\begin{tabularx}{\linewidth}{|X|X|X|}
\hline
\textbf{品牌} & \textbf{模式名稱} & \textbf{打開方法} \tabularnewline
\hline
iPhone & 輔助使用（冇單獨簡易模式，但可逐項調大） & 設定 > 顯示與亮度 > 文字大小，拉到最大。再入「輔助功能 > 顯示與文字大小」開「粗體文字」「按鈕形狀」 \tabularnewline
華為/榮耀 & 簡易模式 & 設定 > 系統同更新 > 簡易模式 > 開啟 \tabularnewline
小米/紅米 & 極簡模式 & 設定 > 更多設定 > 極簡模式 > 開啟 \tabularnewline
OPPO/一加 & 簡易模式 & 設定 > 特色功能 > 簡易模式 \tabularnewline
vivo/iQOO & 簡易桌面 & 設定 > 桌面、鎖屏與桌布 > 簡易桌面 \tabularnewline
三星 & 簡易模式 & 設定 > 顯示 > 簡易模式 \tabularnewline
\hline
\end{tabularx}
\end{center}

如果你用嘅係其他品牌，喺手機「設定」頂部嘅搜尋框入面便直接搜「簡易」或「長輩」兩個字，一般都能搵到。搵唔到嘅話，按下面嘅「手動調大」方案一條條設定，效果都一樣。

\booksection{二、如果唔開簡易模式，手動調都得——「大字+大圖標」五步走}

有啲長輩唔想用簡易模式，覺得介面太「老氣」。咁就手動調以下五項，效果一樣好：

\begin{enumerate}
  \item \textbf{字體調到最大}：設定 > 顯示 > 字體大小，拖到最右邊。順便將「顯示大小」（有啲叫「顯示縮放」）都調到最大，咁樣所有介面元素都跟住變大。
  \item \textbf{圖標變大}：喺主螢幕雙指捏合（好似縮小相片咁樣）進入桌面設定 > 桌面佈局（或圖標大小），選格子最少嘅佈局（比如3×5），圖標自然就大喇。
  \item \textbf{打開「高對比度文字」}：等文字邊緣更鋭利，喺「輔助功能」入面便可以搵到呢個開關。
  \item \textbf{關閉不必要嘅通知}：設定 > 通知，將除咗電話、短訊、微信之外嘅App通知全部關掉。咁樣手機唔會一陣響一下，唔會等 人覺得煩。
  \item \textbf{將常用聯絡人放到桌面}：長按桌面空白處 > 添加小工具 > 聯絡人，將老伴、兒女、醫生、緊急聯絡人嘅快捷撥號放到桌面度。點一下就撥出去，唔使揭通訊錄。
\end{enumerate}

\begin{tcbpractice}{動手：幫長輩「裝修」一下手機}
\begin{enumerate}
  \item 打開簡易模式（按上面嘅品牌表格操作）。
  \item 將電量數字顯示打開（設定 > 電池 > 電量百分比），等佢哋隨時知道仲剩幾多電。
  \item 調出通知欄，將藍牙、NFC、飛行模式呢啲長輩用唔着嘅快捷開關移除或收埋。
  \item 喺桌面放四個快捷撥號：老伴、兒女、家庭醫生（如果有）、社區/物業。
\end{enumerate}

\end{tcbpractice}

\booksection{三、防詐騙——呢章請一定帶住長輩一齊睇}

老年人係電信詐騙嘅重災區。唔係佢哋警惕性低，而係騙子嘅話術經過精心設計，專門針對老年人嘅心理弱點。以下係我整理嘅最常見騙局類型同應對方法，建議你直接讀畀長輩聽，或者列印出嚟貼喺電話隔離。

\booksubsection{六大常見騙局（見到即刻掛電話）：}

\begin{itemize}
  \item \textbf{(1) 「你孫仔出事喇，趕緊匯錢」}——騙子會講「你孫仔/孫女喺出面闖禍喇，被警察拉咗，需要保證金」。記住：\textbf{任何要求轉賬先至可以「放人」嘅電話都係詐騙}。公檢法唔會電話要錢。
  \item \textbf{(2) 「你嘅醫保卡被停用喇」}——短訊入面便講醫保卡異常，等你點連結「激活」。記住：\textbf{醫保局、社保局唔會發短訊等你點連結}。遇到呢種情況直接打12333（社保官方熱線）核實。
  \item \textbf{(3) 「高收益理財，穩賺唔賠」}——微信羣、短訊入面便推薦各種「養老理財」「國家扶持項目」，承諾高額回報。記住：\textbf{承諾年化收益超過5\%仲「保本」嘅，99.9\%係騙局}。
  \item \textbf{(4) 「我係你領導/老朋友，換號碼喇」}——冒充熟人借錢。記住：\textbf{打電話核實比咩都管用}，話「唔方便接電話」嘅八成係騙子。
  \item \textbf{(5) 「免費領雞蛋/保健品講座」}——線下免費活動誘導老人聽「健康講座」，然後賣天價保健品。記住：\textbf{天上唔跌餡餅，免費活動背後都係生意}。
  \item \textbf{(6) 「你嘅手機中毒喇，需要遠程修復」}——騙子等你下載一個遠程控制App（如向日葵、TeamViewer），然後遠程操控你嘅手機轉走錢。記住：\textbf{絕對唔好等任何人遠程操控你嘅手機！}手機廠商、通信運營商、銀行都唔會畀你打電話等你裝遠程控制軟件。
\end{itemize}

\booksubsection{三道防線，立好之後能擋掉90\%嘅詐騙：}

\begin{enumerate}
  \item \textbf{第一道：手機自帶攔截}——安卓手機打開「騷擾攔截」（喺撥號App或手機管家App入面便），iPhone打開「靜音未知來電」（設定 > 電話 > 靜音未知來電）。
  \item \textbf{第二道：運營商攔截}——發送短訊「KTFSR」到10086，免費開通移動嘅高頻騷擾電話攔截；電信聯通亦有類似服務。打運營商客服電話等人工幫你開通就得。
  \item \textbf{第三道：國家反詐中心App}——喺應用商店下載「國家反詐中心」App，註冊並開啟來電預警。呢個App同公安嘅數據庫聯動，騙子號碼一撥過嚟就能識別。
\end{enumerate}

\begin{tcbwarn}{畀長輩嘅「防騙口訣」（列印出嚟貼電話隔籬）}
\begin{itemize}
  \item \textbf{陌生電話要轉賬 掛掉！}
  \item \textbf{短訊連結等你點 刪掉！}
  \item \textbf{話高收益保本嘅 拉黑！}
  \item \textbf{等你裝遠程控制軟件 直接掛！}
  \item \textbf{真警察唔會電話辦案更唔會要錢！}
  \item \textbf{拿唔準嘅事，先畀兒女打個電話問問。}
\end{itemize}

\end{tcbwarn}

\booksection{四、遠程協助——你喺外地都能幫爸媽修手機}

有時候長輩遇到嘅問題，你喺電話入面便描述半日佢哋都搵唔到對應嘅按鈕。呢個時候「遠程協助」功能就派上用場喇——你可以直接喺自己嘅手機上見到佢哋手機嘅螢幕，幫佢哋操作。各品牌嘅方案唔一樣：

\begin{itemize}
  \item \textbf{華為/榮耀}：自帶「暢連」App，入面便有「螢幕共享」功能。你同爸媽都裝暢連，發起視像通話後㩒「共享螢幕」，你就可以喺佢哋螢幕上畫圈標註，話畀你媽媽聽「點呢度」。
  \item \textbf{小米/紅米}：小米視像通話App入面便有「遠程協助」，支援螢幕共享+語音通話同步。
  \item \textbf{OPPO/vivo}：內置「家庭空間」或「遠程守護」，喺設定入面便搜尋「遠程」就能搵到。
  \item \textbf{通用方案（所有品牌都可用）}：下載「向日葵遠程控制」App（免費版夠用），你哋各裝一個。註冊同一個帳號或加好友，等爸媽將手機上顯示嘅驗證碼讀畀你，你就能遠程見到同操作佢哋嘅螢幕喇。
  \item \textbf{iPhone用戶}：FaceTime通話中支援「共享螢幕」。喺FaceTime通話時，㩒螢幕共享按鈕，對方就能見到你嘅螢幕。但iOS安全限制無法遠程操作，淨係可以「指揮式教學」。
\end{itemize}

\booksection{五、微信嘅「適老改造」——同長輩講呢幾句就夠}

微信係長輩哋用得最多嘅App，但介面太密太雜。以下係幾個立竿見影嘅調整：

\begin{itemize}
  \item \textbf{調大微信自己嘅字體}：微信 > 我 > 設定 > 通用 > 字體大小，拖到最大。注意，呢個淨係影響微信內部嘅字體，唔影響手機系統字體。
  \item \textbf{打開「關懷模式」}：微信 > 我 > 設定 > 關懷模式 > 開啟。打開後字體同按鈕都會變大，色彩對比度加強，仲能一鍵「聽文字消息」（點一下文字消息就能朗讀出嚟）。
  \item \textbf{幫佢哋備註好重要聯絡人}：將兒女、親戚嘅名喺通訊錄入面便改成容易認嘅——比如「大仔-小明」「二女-小紅」，等佢哋一打開微信就知邊個係邊個。
  \item \textbf{教佢哋用「語音輸入」而唔係打字}：聊天框隔籬嘅鍵盤上，按住空格鍵（安卓）或點麥克風圖標（iPhone），直接講嘢就能變成文字發送。識別準確率已經好高，對唔識拼音打字嘅老人嚟講係「質變級」嘅便利。
  \item \textbf{關閉「發現頁」入面便唔用嘅功能}：微信 > 我 > 設定 > 通用 > 發現頁管理，將「朋友圈」「影片號」「直播」「購物」「遊戲」入面便唔用嘅全部關掉，減少幹擾。
\end{itemize}

\booksection{六、手機唔見咗點算？幫長輩做個「應急預案」}

預則立，唔預則廢。趁手機仲未唔見，先將呢幾件事做好：

\begin{enumerate}
  \item \textbf{確認「查找我嘅手機」已打開}：iPhone喺設定 > 你嘅名 > 查找 > 查找我嘅iPhone；安卓喺設定搜「查找設備」。務必確認呢個開關係亮住嘅。
  \item \textbf{將重要嘅密碼同帳號記喺一張紙上}：Apple ID / 華為帳號 / 小米帳號嘅密碼，鎖喺屋企抽屜一份，畀信得過嘅兒女一份。唔見咗手機之後要登入呢啲帳號嚟遠程鎖定同定位，如果唔記得密碼就麻煩喇。
  \item \textbf{教長輩一句話}：\textbf{「手機唔見咗，第一時間畀家人打電話。」}家人可以立刻登入iCloud或安卓嘅查找設備網頁，遠程標記為「丟失模式」，咁樣執到嘅人都打唔開、用唔到。
\end{enumerate}

\booksection{七、無線耳機同助聽器相容——被忽視嘅好功能}

好多長輩有聽力問題，但唔知手機可以同助聽器協同工作。iPhone同部分安卓手機支援MFi（Made for iPhone）助聽器，可以直接將電話聲音傳到助聽器入面便，比揸住手機聽筒清楚太多喇。

\begin{itemize}
  \item \textbf{連接方法}：iPhone設定 > 輔助功能 > 助聽設備，打開藍牙，按照助聽器嘅配對說明操作。
  \item \textbf{如果唔用助聽器}：喺「輔助功能 > 音頻/視覺」入面便打開「電話降噪」同「單聲道音頻」，能改善通話清晰度。
  \item \textbf{安卓}：喺「輔助功能 > 聽力增強」入面便有類似設定，各品牌位置略有差異，直接搜「助聽」兩個字。
\end{itemize}

\booksection{八、健康相關功能——手機唔淨係通訊工具}

\begin{itemize}
  \item \textbf{SOS緊急求助}：連續按電源鍵五次（大多數手機），手機會自動撥打緊急電話並畀預設嘅緊急聯絡人發送定位。幫長輩提前設定好——iPhone：「設定 > 緊急SOS」，安卓：「設定 > 安全 > 緊急求助」。教長輩用之前先演習一遍，免得到時緊張按唔準。
  \item \textbf{跌倒檢測}：Apple Watch有呢個功能，檢測到你跌倒且一段時間唔鬱，會自動撥打急救電話同通知家人。有預算嘅話畀長輩配一個，關鍵時刻能救命。
  \item \textbf{用藥提醒}：iPhone嘅「健康」App內置用藥提醒功能，設好藥名、時間同劑量，到點就彈出提醒。安卓可以用「定時鬧鐘」或下載「食藥提醒」類App。
\end{itemize}

\begin{tcbsummary}{本章小結}
\begin{itemize}
  \item 先開簡易模式或手動調大字體/圖標，環境調好大半問題就冇喇
  \item 帶長輩用手機三鐵律：唔好催、用類比、寫低
  \item 六大常見騙局要反覆提醒，三道防線立好能擋掉90\%詐騙
  \item 遠程協助等你人喺外地都能幫爸媽修手機
  \item 微信開「關懷模式」+「語音輸入」係長輩最實用嘅兩個功能
  \item SOS緊急求助同跌倒檢測，關鍵時刻救命嘅配置值得花五分鐘設好
\end{itemize}

\end{tcbsummary}
